% TOP_PROBLEM: {"id": 2690, "title": "Kirby Problem 1.31", "category_id": 7, "category": {"id": 7, "name": "topology", "display_name": "Topology", "description": "Properties preserved under continuous deformations.", "slug": "topology", "order_index": 7, "created_at": "2026-07-31T15:26:25.671Z"}, "status": "open"}
% FIRST_POSTED: null
% ENTRY_KIND: statement_only
% Imported from: batch10.tex; UnsolvedMath ID 2690
\documentclass[11pt]{article}
\usepackage[margin=25mm]{geometry}
\usepackage{fontspec}\setmainfont{Latin Modern Roman}
\usepackage{amsmath,amssymb,mathtools}
\usepackage{unicode-math}\setmathfont{Latin Modern Math}
\usepackage{xurl,hyperref}
\hypersetup{colorlinks=true,urlcolor=blue}
\setcounter{secnumdepth}{0}
\setlength{\parindent}{0pt}\setlength{\parskip}{5pt}
\setlength{\emergencystretch}{3em}
\newcommand{\sref}[2]{\hyperlink{source:#1:#2}{[S#2]}}
\newcommand{\cataloguescope}[1]{\paragraph{Recorded target.}#1}
\begin{document}
% UnsolvedMath ID 2690; source catalogue code KP-1.31
\setcounter{section}{2689}
\section{Kirby Problem 1.31}\label{2690}
{\small\sffamily UnsolvedMath ID 2690 · Topology}
\subsection{Definitions and mathematical statement}

\paragraph{Shared notation.}$\mathbb N=\{1,2,\ldots\}$, and $\mathbb Z,\mathbb Q,\mathbb R,\mathbb C$ denote the integers, rationals, reals and complex numbers. Any other conventions are specified locally.


All knots and links below are smooth in the oriented $3$-sphere, up to ambient isotopy, unless another ambient manifold is specified. For coprime nonzero integers $m,n$, $T(m,n)$ is the torus knot of those parameters; changing the sign of one parameter gives its mirror, and a parameter of absolute value one gives the unknot. Write $\mathrm{Kh}^{i,j}(K;\mathbb Z)$ for integral Khovanov homology, with homological and quantum gradings. Write $\mathrm{KhR}_{N}^{i,j}(K;\mathbb Z)$ for integral fundamental-representation $\mathfrak{gl}(N)$ Khovanov--Rozansky link homology, for $N\geq2$, using the integral link-homology constructions discussed in the source. A computation of these groups includes all graded free ranks and torsion summands, rather than only their Euler characteristics.
\paragraph{Statement recorded in the supplied source.}
Determine (a) $\mathrm{Kh}^{i,j}(T(m,n);\mathbb Z)$ for every torus knot and all $i,j$, and (b) the corresponding $\mathfrak{gl}(N)$ Khovanov--Rozansky homologies for all torus knots and all $N\geq2$, including the torsion phenomena emphasized in the original remarks. \sref{2690}{1}\sref{2690}{2}
\subsection{Short English statement}
Compute all graded homology groups, including integral torsion, for every torus knot in the Khovanov and general linear Khovanov--Rozansky theories.
\subsection{Sources}
\begin{description}
\item[\hypertarget{source:2690:1}{S1}] ULAM.AI and the UnsolvedMath Contributors, \emph{UnsolvedMath: A Curated Collection of Open Mathematics Problems}, record 2690 (KP-1.31). CC BY 4.0. \url{https://huggingface.co/datasets/ulamai/UnsolvedMath}.
\item[\hypertarget{source:2690:2}{S2}] R. \.{I}nan\c{c} Baykur, Robion C. Kirby and Daniel Ruberman, K3 -- A New Problem List in Low-Dimensional Topology, author preliminary edition (2026), Problem 1.31 and its remarks. \url{https://math.berkeley.edu/sites/default/files/surv-295-ruberman-watermarked-author-pdf.pdf}.
\end{description}
\label{end:2690}
\end{document}
